\section{Limitations and next steps}
\label{sec:limitations}

The paper makes bounded claims:

\begin{itemize}
  \item \textbf{Control.} \emph{Scope:} The exact results assume affine
  readouts, one fixed direction, proportional control, and no clipping. The
  nonlinear test reruns a learned Transformer block, but the model is small and
  its latent features are supervised; the 29.0\% gain belongs to the matched
  estimator--direction pair, not the estimator alone. \emph{Next test:} extend
  the comparison to token-level control in larger models with independently
  usable edit directions.

  \item \textbf{Pretrained-model interventions.} \emph{Scope:} The GPT-2 and
  Qwen studies use pinned models, selected head sets, fixed ablations, and narrow
  prompt distributions. The IOI confirmation is pilot-informed and locally
  sealed rather than independently preregistered. The Qwen task keeps only
  eligible prompts; full-vocabulary top-1 accuracy on unfiltered prompts is
  24.45\%, and the direct-loss observer ties the natural-mean observer at the
  largest target.
  \emph{Next test:} use more models, unseen prompt templates, and independently
  fixed candidate sets and targets.

  \item \textbf{Authorization safety.} \emph{Scope:} The inert task asks Qwen
  directly for an authorization verdict, so it cannot establish a uniquely
  hidden safety state or test harmful execution, deception, or adaptive attacks.
  Its held-out failure also shows that one global threshold can hide poor
  calibration between groups. \emph{Next test:} remove the explicit verdict
  question and provide calibration data for every intended deployment group
  while keeping the model, budgets, and evaluator fixed.

  \item \textbf{APPS backdoors.} \emph{Scope:} The study uses stored working
  backdoors rather than attacks adapted to each monitor. One model supplies all
  learned representations, the text baseline is simple, and the prevalence
  sweep keeps the dataset-family mix fixed. \emph{Next test:} freeze generated
  attack trajectories and training budgets, then add adaptive evasion, stronger
  code monitors, and a deployment-level shutdown action.

  \item \textbf{Artifact.} \emph{Scope:} The release provides checked control,
  finite-effect, and safety interfaces, but it is not yet a general task
  ecosystem. \emph{Next test:} add literature baselines, another safety
  environment with adaptive attacks, and a stable end-to-end reproduction path.
\end{itemize}

\section{Conclusion}

ObserverBench asks a simple question: is an estimate good enough for the action
chosen from it? An observer need not reproduce every detail of a model. It must
preserve the information that the available action, decision rule, and loss
actually use.

The experiments show why this distinction matters. A better effect predictor
can choose no better intervention when it predicts the wrong quantity. A
perfect classifier can spend a safety budget poorly when it ignores consequence
or calibration between groups. A simple output signal can also make an internal
probe unnecessary. ObserverBench records these choices in versioned tasks and
applies the same fixed policy to every submission; it does not rank observers in
the abstract. Before trusting an observer, state what it may read, freeze the
action rule and budget, and score the outcome the action actually produces.
